% Version 3.10: finite opening selection before Ready; no live-pool REPORT scan.
\subsubsection{Prepared opening input}
\label{sec:input-selection}
\label{sec:prepared-input}
Let $A=A_{\beta,r,0,0}$ be the exact effective DATA opening anchor,
and let $N=V_i(r)$ be the target snapshot authenticated by
$\tagg{AUTH}_{r,0}$. Use its original weights and fault bound $F$.
Order canonical authenticated requests by
\[
 \operatorname{key}(x)=
 (|\operatorname{enc}(x)|,\Hash(\operatorname{enc}(x))),
\]
lexicographically, with $\tagg{EMPTY}$ after every request key.
Hash binding identifies the canonical body. Each fixed request has
finitely many predecessors, because their byte strings have bounded
length. This gives no practical request-latency bound.
A request is \defterm{eligible at $A$} if its full authenticated bytes
pass the original application admission check immediately after $A$,
it fits alone in DATA, and its logical identity is absent from the
complete ancestry of $A$.

The immutable preparation key is
\[
\begin{split}
 K={}&(\mathit{chain},\mathit{protocol},\tagg{PREPARED\_INPUT},i,r,
       \Hash(\tagg{AUTH}_{r,0}),\\
     &A.\mathit{original\_pair},\Hash(A),h(A),k(A),
       \mathit{native\_validation\_domain}).
\end{split}
\]
The last field names the existing native validation interface; it does
not invent a state commitment. The original evidence connects the actual
executable object to $A$. The key excludes the target branch version,
recovery, term, QC encoding, receipt time and private queue root. The
original producing pair of $A$ remains part of its identity. Reuse
requires equality of every displayed field; a different exact anchor,
authority or native domain starts different work.

\paragraph{Durable preparation state.}
The request owner of \cref{sec:request-ownership} retains
\tagg{RawJournal}, canonical requests and their original first-receipt
metadata, pending entries, and the full checked hydration certificate.
A complete received raw record has an immutable stored-byte handle,
length, digest and authentication material. Fragment completion appends
that handle before releasing its receive obligation, even if ordinary
canonicalization has not run. This is the ownership journal's incremental
receipt hook, not application admission or a whole-pool scan.
Canonical records already retained remain available through a versioned
canonical-store cut. A globally unaccepted request therefore has a raw
or canonical handle even when the public-import worker is behind.

For each $K$ and member $v$, the same owner retains a verified executable
state \tagg{VerifiedState}[$K$]; at most one finite
\tagg{SelectionJob}[$K,v$] with its cut, cursor, best candidate and saved
work state; at most one signed \tagg{Offer}[$K,v$]; and any complete valid
\tagg{OpeningCert}[$K$]. An initial Enter also retains a session attachment
containing its executable state and one matching certificate.
Worker exit, recovery and branch retry preserve these records and their
service duties. Verified inclusion in $G$ may remove a pending entry,
but not bytes pinned by a job, evidence, certificate or outgoing duty.

\paragraph{State distribution before any Ready.}
As soon as a correct preparer has constructed and validated the original
executable $A$-state, it actually multicasts
\[
 \tagg{STATE\_CANDIDATE}
   (K,\text{original complete Ready state and dependencies})
\]
to the target members. The packet contains the executable object and
all original exact-anchor, native, authority, handoff and source evidence.
It has no opening certificate. The receiver validates those original
dependencies and may cache the executable object as
\tagg{VerifiedState}[$K$]. It grants no entry, STATUS, ANCHOR vote,
progress timer or lock; public evidence still reaches its original
manager handler, including closure processing.

Every correct native preparer sends without waiting for an offer,
certificate, Ready or acknowledgement from another preparer. Candidate
construction, actual validation, any required native conversion and any
finite fragmented transfer occur before Ready, with no new $\Delta$
deadline. A received executable object can supply the state dependency
of selection without an independent reconstruction. Thus selection
does not wait circularly for the first ReadyBundle. The original eventual
state-availability premise is still needed for a correct preparation to
finish. This eventual preparation path does not establish the separate
timed direct-adoption contract of \cref{sec:ready-bundle}.

\paragraph{A fixed finite selection job.}
For $K$, correct $v$ performs the following persistent procedure.
\begin{enumerate}[label=(\roman*),itemsep=2pt]
\item Authenticate the target and anchor, and start or receive the
original executable-state preparation concurrently with hydration.
After fully checking a matching \tagg{HYDRATED} certificate, freeze the
current finite \tagg{RawJournal} prefix and canonical-store version.
Use maintained counters/version handles to take the cut; do not run a
scan in the certificate-receive handler. The cut includes every complete
retained raw record already received and every retained canonical entry,
including records awaiting ordinary canonicalization. Later arrivals
never enlarge or restart it.
\item Once \tagg{VerifiedState}[$K$] is actually available, examine every
record of the frozen cut. Check canonical encoding and authentication,
merge only validated duplicate identities without changing FIFO metadata,
and persist new valid entries through the same first-authenticated-receipt
and individual-relay hook of \cref{sec:request-ownership}. Thus new
canonical pending entries also enter its \tagg{ReceiptJournal}.
Evaluate eligibility at the exact $A$ using
its verified state and complete ancestry. Any ancestry-index construction,
record assembly, admission work and minimum search belong to this actual
preparation. Skip malformed or ineligible records without discarding
still-pending requests. No missing body or future arrival is awaited.
\item Finish the finite job, persist \tagg{SelectionDone}[$K,v$] and its
result, and sign exactly one
\tagg{OPENING\_OFFER}$(K,k_v)$, where $k_v$ is its least eligible key,
or \tagg{EMPTY} if none exists. A nonempty offer supplies the full
canonical authenticated body. Actually multicast this complete
contribution to all target members; unfinished sends remain durable.
A correct member never signs before completing the job.
\end{enumerate}
Each job has a fixed finite input and terminating local checks. The
owner continues its saved finite steps until completion, independently
of new arrivals and worker changes. A large finite input can take
arbitrarily long; no uniform CPU, storage or preparation bound follows.
Completion certifies only this cut at $A$, not a live pool or another
parent's eligible projection.

\begin{definition}[Opening certificate]
\label{def:opening-certificate}
A complete opening certificate has the form
\[
 \mathit{OC}=(K,\{(v,\langle\tagg{OPENING\_OFFER},K,k_v\rangle_v)\}_{v\in S},
                x_{\rm selected}).
\]
Its distinct target members have total weight $>F$. If any supplied key
is nonempty, $x_{\rm selected}$ is the complete authenticated eligible
body with the minimum supplied key; otherwise every key is \tagg{EMPTY}
and no selected body is present. The common predicate
$\tagg{CheckOC}(A,\tagg{AUTH}_{r,0},\mathit{OC})$ checks the exact $K$,
membership, distinctness, signatures, weight, minimum of the supplied
keys and the selected body's canonical key and eligibility at $A$.
It does not read a private queue, obtain nonselected bodies or rerun
any signer's selection job.
\end{definition}

Each correct collector checks complete offers, including each nonempty
contribution's body and eligibility, and counts a member once.
Equivocation cannot increase its weight or restart the collector.
As soon as retained valid contributions exceed $F$, form, persist and
actually relay the complete certificate to all target members. Finite
fragmentation may occur before Ready; it does not turn the whole
transfer into one $\Delta$. Every recipient verifies and retains the
complete certificate. A correct member may adopt any valid matching
certificate and cancel redundant unfinished selection, while preserving
raw/pending data and all promised storage, serving and relay duties.
Absent an already valid certificate, it continues its finite job and
offer dispatch.

The Ready input gate is this complete $\mathit{OC}$, not a second
\tagg{HYDRATED} signature set. Every valid $\mathit{OC}$ contains a correct
offer signer, and that signer first checked hydration and finished
selection. Its causal history supplies this assertion; cryptographic
inspection alone does not prove a minimum over an unknown pool. The
receiver need not obtain the signer's cut, hydration certificate or
historical directory. A locally attached certificate is not exclusive:
a later valid NEWVIEW or proposal can use a different valid certificate
for the same $K$ without a private-certificate veto.

\begin{proposition}[Finite opening preparation]
\label{prop:opening-formation}
For a fixed target and exact $K$, once the ownership preparation of
\cref{prop:ownership-formation} and the original eventual executable-state
preparation complete, a complete opening certificate is eventually
available at every correct target member. A retry with the same $K$
does not repeat that work.
\end{proposition}
\begin{proof}
A correct native preparer distributes its complete candidate before
waiting for offers. Each continuing correct selection job then freezes
and processes one finite cut. Unless a certificate already exists,
every correct member finishes and actually sends its offer. There are
finitely many target members, so their finite completion and delivery
times have a finite maximum. Correct weight is
$W-B\ge W-F>F$; every correct collector can form a certificate without
Byzantine cooperation or knowing an ``all correct replies'' event.
Early adoption can only make a certificate available sooner. The key,
saved cursor, state, offers and certificate survive branch retries.
\end{proof}

\begin{proposition}[Opening coverage]
\label{prop:opening-coverage}
Fix an unaccepted request $x$ and a nonexceptional target in
\cref{prop:ownership-coverage}. If $x$ is eligible at its accepted opening
anchor $A$, every valid opening certificate for $K$ is nonempty and
selects a request $y$ with
$\operatorname{key}(y)\le\operatorname{key}(x)$.
\end{proposition}
\begin{proof}
Every valid hydration certificate for this target contains a correct
signer whose completed dissemination puts $x$'s complete bytes and
raw handle at every correct target before that certificate exists
(\cref{prop:ownership-coverage}). Every correct offer freezes its cut
after checking such a certificate. While $x$ is absent from $G$, its
raw or canonical handle is therefore in that cut, and the completed
eligible-at-$A$ search selects $k_v\le\operatorname{key}(x)$.
Every valid $\mathit{OC}$ contains a correct offer because its weight
exceeds $F$. Its minimum is at most that correct key, and
\tagg{CheckOC} supplies an actual eligible body for the minimum.
No honest intersection between hydration, offer and report signer
sets is used; actual all-target dissemination connects them.
\end{proof}

In a fixed-$G$, fixed-slot tail, hydration, candidate preparation,
selection and certificate formation add finite completion times to
the preparation maximum. They add no post-Enter pool-processing term
to the regional response bound. The same $K$ survives zero-output
branch retries; fresh view-bound advertisements and reports only read
the certificate. A speculative recovery parent either retains its
first own-session DATA witness or returns to $A$ and reuses the same
certificate. No selection action resets an export window, Force cutoff,
entry time or progress timer.
